\appendix

\chapter{物理原型实现与 baseline 调优审计}
\label{app:physical-implementation-audit}

\parab{CPU 资源口径.}
每次正式运行记录 rank、角色、poller 数、EAL lcore、NUMA socket、NIC 队列和
内存放置。正文以执行协议数据面的已配置 poller lcore 为资源单位；只负责
发布 collective 和等待完成的 EAL main lcore 不计入该指标。四台主机在运行前
固定 performance governor，并对实际选用的 lcore 逐核核对；不满足 governor、
NUMA、进程互斥或链路门禁的运行不进入正式统计。我们不使用进程生命周期上的
CPU-time 比值，因为初始化、测量和退出阶段的分母并不对应同一数据面窗口。

\parab{镜像绑定与两层参数.}
正式绑定为：Arbor 单层结果使用 r81 镜像，细粒度物理状态扫描与其后的
端侧 fast-path 实验使用通过同等 route/timing 门的 r88；ATP-8K 与
SwitchML-8K 分别使用独立的 r87 和 r79；两层实验使用 r60。两层实验固定
Swift target 40~\textmu s、$W_0=2$、cap 16、8-\textmu s startup
stagger、12-credit shared window 和 200-\textmu s stall repair；端点
日志中的 \texttt{slot\_pool\_size=1024} 只是可动态映射的逻辑标识
空间，不计作物理 row。归档 fault-proxy 批次的 35 个成功 attempt 之外，
另有三个 attempt 在发包前因基础设施 preflight 被拒，不计入成功统计。

\parab{Baseline 优化边界.}
ATP-8K 与 SwitchML-8K 使用各自独立的端侧协议状态机和 FPGA 镜像。端侧候选
包括 borrowed RX、NUMA-local payload、每 poller 独立队列、批处理、状态查找
快速路径和多段 TX；每个候选先通过 exact-FP32、RX/TX 守恒和 poller 所有权
检查，再用独立 engineering pilot 判断是否保留。无收益或降低吞吐的候选仍
保留源码、命令和原始日志，不从审计中删除。Arbor 的 credit 调度、repair
触发和共享 allocator 不移植到 baseline，因而端侧优化不会改变被比较算法的
恢复或资源分配语义。

正式 CPU 扫描保留 Arbor 和 SwitchML 的 1/2/4/8-poller 曲线，以及 ATP
的 1/1、2/2、4/4、8/4 和 8/8 PS/worker 配置。每个点绑定源码、binary、
PDI/UUID、lcore、NUMA 策略、CPU governor 和五次原始吞吐；正文只从该
冻结曲线选择首个达到 90~Gbps 的配置。
SwitchML 的 50 次正式端点运行均通过 exact-FP32；整批上行共恢复
322/18,032,518 个包（17.86~ppm），下行丢包、Tofino TM drop 和端点硬错误
均为零。因此正文把该批次标为 recovery-qualified，而不写成零丢包。

\parab{SwitchML 端侧关键路径审计.}
保留实现让 RX result payload 在 callback 生命周期内直接借用 mbuf，并把
slot 查找、generation 校验、完成迁移和 replacement 生成合并为单写者
\texttt{ConsumeResult} 快路径；scheduler 也长期持有连续输入区，避免逐包
共享指针增减。我们另外测试了显式 \texttt{memcpy}、AVX-512、poller 与内存
全部放置在 NIC NUMA~0、non-temporal store、强制 TX coalescing，以及外部
payload 的两段式 zero-copy TX。前几项没有增益或降低吞吐；两段式 TX 虽
消除了 payload copy，却增加 mlx5 multi-segment WQE 与 TX burst 开销，因而
不进入正式实现。保留快路径的 p4 profiling 中，result callback 与线性 TX
serialization 合计约 5.2--5.7k cycles/packet。由于 SwitchML 每个 slot 的
下一份贡献由当前 result 完成后释放，这些工作直接位于 slot 的自时钟进度
路径；p8 通过并行独立 slot 周转达到线速。该审计解释本原型的核数曲线，
不把 p8 表述为 SwitchML 算法或其它 NIC 实现的固有下界。

\parab{最小 poller 配置下的原语覆盖.}
正文的端侧扫描选择 Arbor R2/P2，即每 rank 两个 datapath poller、两个
channel replica，全系统共八个 poller。为检验这一最小配置是否只对
AllReduce 有利，我们固定四个 rank、单 job 和 $A=1024$，对五种原语与
六种消息大小运行完整笛卡尔积。warm-up 使用独立的 lane~0，正式迭代
轮转访问每 rank 预注册的 128-MiB buffer ring。每种原语在正式矩阵前
冻结传输参数；每个 cell 随后运行五个新进程，不按消息大小重新调参。
端点可按 primitive 和 message size 选择 CPU-copy 或同 NUMA I/OAT
delivery path；16-MiB AllGather 及 64-MiB AllGather/AllReduce 的选择由
独立交替配对实验冻结，不增加 poller。Reduce
的 P2 16/64-MiB 点使用等字节 cold-payload 对照：前者连续执行四次
16-MiB Reduce 并访问四个不同区域，后者执行一次连续 64-MiB Reduce；
两者每 rank 均测量 64~MiB，协议 warm-up 使用另一块内存。

\begin{table*}[t]
\centering
\small
\caption{Arbor 在最小 R2/P2 端侧配置下的五原语与消息大小矩阵。每格为
五个新进程的中位 logical goodput（Gbps）。AllGather 和 ReduceScatter
报告 API goodput / $N=4$ 瓶颈链路等效 goodput；其余原语的换算系数为 1。}
\label{tab:physical-workload}
\setlength{\tabcolsep}{3.5pt}
\resizebox{\linewidth}{!}{%
\begin{tabular}{lrrrrrr}
\toprule
原语 & 64~KiB & 256~KiB & 1~MiB & 4~MiB & 16~MiB & 64~MiB \\
\midrule
% Generated by analyze_cache_diverse_transport_rerun.py, with the P2 Reduce
% 16/64-MiB cold-payload override recorded in the installed provenance.
AllReduce & 7.489 & 23.622 & 50.890 & 77.749 & 90.359 & 97.217 \\
Reduce & 10.303 & 28.479 & 57.202 & 73.295 & 85.109 & 86.592 \\
Broadcast & 11.901 & 33.514 & 64.950 & 85.838 & 93.564 & 96.580 \\
AllGather & 8.136 / 24.407 & 16.593 / 49.778 & 23.720 / 71.160 & 28.213 / 84.640 & 30.154 / 90.462 & 30.851 / 92.553 \\
ReduceScatter & 7.547 / 22.642 & 16.925 / 50.775 & 26.605 / 79.816 & 31.055 / 93.166 & 32.396 / 97.189 & 32.761 / 98.283 \\
\bottomrule
\end{tabular}}
\end{table*}

AllGather 和 ReduceScatter 的 API 分子只计一次用户消息；四-rank 下每个
rank 的瓶颈方向需要搬运三份 peer 数据，故链路等效值为 API goodput 的
三倍。跨原语比较使用斜线右侧：64-MiB AllGather 的 30.851~Gbps API
goodput 对应 92.553~Gbps 链路等效 goodput。64-MiB 的交替配对中，I/OAT
相对 CPU-copy 将 AllGather/AllReduce 的中位数从 89.139/93.969 提高到
92.553/97.217~Gbps，两者均为 5/5 胜。所有正式样本均保留恢复事件，
并通过数值与 DMA 门；物理计数器缺口
须能精确归因并在表后披露，否则样本失效。表中数值不表示零恢复事件。
除上述两个 Reduce cell 外，该口径是
cache-diverse steady state：lane~0 只预热协议，正式阶段遍历预注册的
128-MiB ring。等字节对照中，四个不同 16-MiB 区域和连续 64-MiB 区域
分别达到 85.109 和 86.592~Gbps；独立五次确认得到 84.970 和
86.069~Gbps。重复同一 16-MiB 区域达到 90.044~Gbps，但该 payload-reuse
对照不进入表格。两批确认均逐次通过数值校验；硬件封口分别记录
4/184,320 和 2/184,320 帧的 converter ingress 缺口，故这两个 cell 为
recovery-qualified，而非物理零丢包。

Reduce 的大消息路径还执行一次只改变 responder 并行度的 P2/P4 对照：

\begin{table*}[t]
\centering
\small
\caption{Arbor Reduce 的 responder poller 对照。P2 的 16/64-MiB cell
使用等字节 cold-payload 对照；其余 cell 使用 cache-diverse buffer ring。
每格为五次运行的中位 logical goodput（Gbps）。}
\label{tab:physical-reduce-pollers}
\setlength{\tabcolsep}{3.5pt}
\resizebox{\linewidth}{!}{%
\begin{tabular}{lrrrrrr}
\toprule
Responder poller & 64~KiB & 256~KiB & 1~MiB & 4~MiB & 16~MiB & 64~MiB \\
\midrule
2/rank & 10.303 & 28.479 & 57.202 & 73.295 & 85.109 & 86.592 \\
4/rank & 9.005 & 28.804 & 60.418 & 85.205 & 95.155 & 97.869 \\
\bottomrule
\end{tabular}}
\end{table*}

R2/P2 在 16 和 64~MiB 时分别达到 85.109 和 86.592~Gbps，未出现随
消息增大而下降。只把 Reduce responder 扩展到每 rank 四个 poller 后，
16 和 64~MiB 分别达到 95.155 和 97.869~Gbps；镜像、wire protocol 和
拥塞参数均不变，因此差值来自 endpoint reduction service capacity。
该对照不改变正文的跨系统 AllReduce 和最小 poller 结论。64~KiB 的 P4
结果从 10.303 降至 9.005~Gbps，增加 poller 并非所有消息大小都受益。

% TODO(testbed-v80-resource-appendix): 在最终镜像与硬件运行封存后加入完整
% Vivado hierarchy 表。分别报告 deployed whole image、公共 shell、DCMAC/
% converter buffering、payload accumulator 与协议控制 dataplane 的
% LUT logic/LUTRAM/FF/BRAM18/BRAM36/URAM/DSP；绑定 PDI、routed DCP、UUID、
% 原生 timing、NoC 和 source hash。r74 的 2048-depth converter 与每 bank
% 512-entry completion-descriptor FIFO 必须单列，不能归因于 accumulator。

\chapter{仿真 baseline 模型与配置校准}
\label{app:sim-models}

\parab{ATP-like.}
ATP-like 以 \texttt{CRC32(job,index)} 映射逻辑聚合表；碰撞后，PS 从尚未
声明的位置中选择新位置，无法在交换机完成的输入由端侧 reducer 处理。
多机架下只有 worker ToR 和 PS ToR 执行聚合，上行 data packet 使用
ECMP，parameter packet 沿固定树返回。该模型保留 ATP~\cite{atp} 的动态
remap、端侧 fallback 和恢复语义，RTO 或连续三个高序 parameter packet
触发当前缺失区间的重传；它不另设小于物理 slot pool 的
congestion-window 上限，也不限制独立逻辑 hash 表或端侧 reducer 吞吐。

\parab{SwitchML-like 与 p2p-CCL.}
SwitchML-like 为每个 job 构造一棵固定的分层聚合树，树上的 ToR、中间层
和 root 均执行聚合。p2p-CCL 采用四条 NCCL-style ring，相邻 channel
方向相反；每条有向 ring edge 保留一条 persistent RDMA queue pair
（QP），并传输 reduce-scatter 与 all-gather 合计 $2(N-1)/N$ 倍的
channel payload。该模型使用 4096-B RDMA payload MTU，并在连续
collective 间保留 QP 与 DCQCN 状态；各 workload family 在独立
calibration 中冻结 DCQCN 参数。

\parab{拓扑深度比较.}
三个 INC 模型均使用 1536 个 8-KiB 物理 slots。\sysname 在两/三层为每个
job 固定四棵候选树；SwitchML-like 每个 job 固定一棵分层树，并把八个 job
按每个 root 两个 job 的方式均衡分配。ATP-like 由组内一个 rank 兼任 PS，
不增加 endpoint；worker/PS ToR 执行聚合，上层 data packet 使用 ECMP。

\parab{LLM training baseline 的拥塞控制校准.}
NCCL 及 ATP/SwitchML hybrid 的 RDMA fallback 使用 packet-level DCQCN。
我们在正式结果生成前声明有限且可复现的调参预算：三个 DCQCN profile 分别
覆盖 AllReduce 与 AG/RS、flat 与 hierarchical NCCL，并各运行三个独立
\texttt{RngRun}，共 54 个 calibration 点。每个 workload family 在通过
DAG、字节、QP 和队列完整性检查的候选中，选择几何平均 goodput 最高的
profile；随后在 $J=1,8,16$ 和新的 \texttt{RngRun=4} 上运行 18 个 held-out
点，再冻结正式配置。正式实验不根据结果重新调参。

选参分数不作为带宽结果报告。冻结的 flat NCCL 在 128 ranks 上取得
188.688~Gbps AllReduce goodput，为其
$400\times128/(2\times127)=201.575$~Gbps ring 上限的 93.607\%；
AllGather 和 ReduceScatter 分别取得 358.857 和 376.879~Gbps，为各自
$400\times128/127=403.150$~Gbps ring 上限的 89.013\% 和 93.484\%。

运行时审计同时核对配置和控制器遥测：fallback QP 必须报告
\texttt{CcMode=1} 及冻结的 DCQCN 参数。Classic 3D direct-tail 实验中，
所有 P2P QP 的最低测量期平均速率达到 400-Gbps 初始速率的 95\% 即记为
无实质降速；瞬时最小速率、ECN 标记、PFC 事件和最终队列排空仍分别报告。
该方法得到的是在公开预算下经独立验证的部署配置，不声称对 DCQCN 参数空间
取得全局最优。

\parab{Classic 3D 的 host-NIC 调度约定.}
\label{app:classic-pp-scheduler}
四种系统执行相同的 PP ready event、payload、route 和训练依赖。Arbor 的
host-NIC egress 另采用一个显式、work-conserving 的次序：ACK/CNP control
仍具有最高优先级；当前可发送的 PP RDMA QP 先于已排队的 Arbor data；PP
不存在、被 DCQCN pacing、window 或 PFC 阻塞时，Arbor 立即使用链路。
Arbor 与 PP 在交换机中仍使用相同的 PG，该约定不预留带宽，也不改变交换机
queue、ECN、PFC、routing 或 collective DAG。

该约定来自 U4/J1 的单变量 canary。旧 host queue 把 Arbor data 放入
ACK-priority queue，持续的 multipath AllReduce 因而可在源 NIC 上把 PP
RDMA 延迟数十毫秒；canary 只打开
\texttt{pp\_preempts\_arbor=1}，随后五个正式 Arbor cell 均冻结相同设置。
正文将它视为训练组合器的实现约定，而不是动态 SRAM 共享机制；用于方法选择
的 U4/J1 canary 被按 exact-run 身份提升，因此不计作独立 held-out 重复。

\parab{50\%-active 轨迹与停止条件.}
八 job 轨迹枚举 $\mathrm{GF}(2)^3$ 上全部 14 个 affine hyperplane：每个
epoch 含四个 job，每个 job 出现七次，任意 job 对共同出现三次。因此任意
均衡的 SwitchML-like 两-job-per-root 配对具有相同的累计共活跃次数，正文
映射不是从性能结果中选择。活动事件只关闭下一轮 collective 的 launch
gate；已启动 collective 继续排空。192~ms 的最后一个事件关闭全部 gate，
随后不再启动配置 backlog 中的后继。审计要求每个 job 完成从 iteration 0
开始的连续前缀、所有已启动 collective 完整提交一次、未启动项无协议状态，
且最终队列和聚合状态排空。

\chapter{在线选树的边界与补充结果}
\label{app:multipath-details}

\parab{时间尺度边界.}
我们还运行了 1、4、16 和 64 个 $R_{fb}$ 的对称轮换。由于每个
256-MiB collective 的无干预完成时间约为 11.3~ms，nonblocking static 的
各树队列可以先发放其它树上已经分配的未来 offset；这些短事件中 uniform 与
dynamic 的 completed-offset goodput 相同。正文采用的 $512R_{fb}$ 是覆盖已审计
隔离完成时间的最小二次幂周期，而不是从旋转实验结果中挑选。这个结果限定了
机制收益：反馈选择减少持续到 collective 尾部的路径失配；更短且能被已有
offset backlog 吸收的瞬态本来就不应显示收益。

\parab{背景流量与方向性锚点.}
64 条与 collective endpoint 不重叠的 DCQCN 长流从 0~ms 开始，并在
1.6--4.8~ms 的共同窗口按每条 QP acknowledged bytes 计量。单 job 的
10\%、25\% 和 50\% offer 下，背景 aggregate goodput 分别为 1.245、3.101 和
5.833~Tbps，dynamic 与 uniform 完全相同。八个 job、25\% offer 时，两者均为
1.141~Tbps，per-flow Jain 均值仅 0.370；因此证据只表明 dynamic 没有造成
额外伤害，不能解释为背景流量不受 collective 竞争。全部 160 个点均无丢包或
PFC。ReduceScatter 和 AllGather 的单 job 锚点中，两种策略都约为
200~Gbps；fixed-window counter 的边界量化值分别为 201.1 和 202.3~Gbps。
方向变化没有引入回退，也没有显示选树收益。

\chapter{共享循环 allocator：记账约定、包络条件与端口界}
\label{app:allocator-port-bound}

\parab{记账约定.}
\Cref{sssec:slot-bound} 的精确条件依赖以下约定。保护区间长度上界
$R_e$ 至少包含一个 allocator arbitration quantum。allocation 与各
Context 的 $t_f$ 事件在交换机 $e$ 内形成全序；同一时刻先处理 $t_f$
事件，再处理复用 allocation。无法保证该同周期顺序的实现应在 $R_e$ 中
再包含一个 arbitration quantum。若 repair cutover 早于 slot 绑定，则
取 $t_f=t_a$，保护区间为空。$N_e$ 对保护区间的计数从 $x$ 自身的
allocation（含）到其 $t_f$（不含）。

\parab{推进包络给出另一条充分条件.}
若实现或 shaper 对任意区间 $I$ 保证
\begin{equation}
 N_e(I)\leq\left\lceil B_e+\lambda_e|I|\right\rceil ,
 \label{eq:allocator-arrival-envelope}
\end{equation}
其中 $B_e$ 是 advance burst，$\lambda_e$ 覆盖该 allocator 的全部
normal advance，则
\begin{equation}
 A_e\geq\left\lceil B_e+\lambda_eR_e\right\rceil
 \label{eq:allocator-depth-envelope}
\end{equation}
满足 \cref{eq:allocator-exact-condition}。$\lambda_e$ 必须是实现保证
的严格包络，而非 trace 平均速率。该条件不依赖正文的 payload 一一
配对前提；在包络保证仍成立的前提下，同样可以代入部署契约保证的
$R_{\max}$。

\parab{端口界的证明.}
对 fixed-$P$ normal-data workload 中的每个 context $y$，在交换机 $e$ 选择
一份与其一一对应、至少携带 $P$ bytes 的 packet $q_e(y)$。
\texttt{Broadcast} 和 \texttt{AllGather} 可选择携带 payload 的 response；
\texttt{Reduce} 和 \texttt{ReduceScatter} 选择 credit 授权的 request；
\texttt{AllReduce} 可选择其中一个方向。每份 packet 只按其实际经过的一条有向
端口计数，multicast copies 不重复计数。记 $t_q(e,y)$ 为该 packet 的末 bit
完成所选有向端口串行化的时刻；同一端口上相邻这类事件至少相隔一个 packet 的
串行化时间。要求
\begin{equation}
 t_a(e,y)\leq t_q(e,y)<t_a(e,y)+R_e.
 \label{eq:allocator-payload-witness}
\end{equation}

固定任一 context $x$。在 $[t_a(e,x),t_f(e,x))$ 内分配的 context $y$ 满足
$t_a(e,y)<t_a(e,x)+R_e$，其配对 packet 又满足
$t_q(e,y)<t_a(e,y)+R_e$。这些 packet 因而都经过长度为 $2R_e$ 的区间
$[t_a(e,x),t_a(e,x)+2R_e)$。端口串行化限制该区间内能够经过的 fixed-$P$
packet 数，逐端口计数即得到 \cref{eq:allocator-directed-port}。若推进
allocator 的 packet 本身携带 payload，则 $t_q(e,y)=t_a(e,y)$，全部配对
packet 已位于长度为 $R_e$ 的 allocation window 内，从而得到
\cref{eq:allocator-payload-carrier-bound}。

\chapter{拥塞控制模块接口}
\label{app:rate-modules}

\Cref{ssec:cc} 为每个 subchannel 提供速率门 $r_s$ 和窗口门 $W_s$。端侧拥塞控制模块可以
更新其中一个或同时更新二者；未更新的速率门取链路线速，窗口门取配置的安全
上限。该模块读取 credit 发出与完成时间、发出时与当前的在途量，以及 repair
事件。基于 completion delay 的算法只使用 responder 本地状态；ECN 算法还接收
各分支 CE 的 OR。控制律不改变 credit、聚合位置或恢复状态。

DCQCN-style group ECN、TIMELY、Swift 和 Oscar 可作为该接口的四个实例。前者
使用显式 ECN，其余三个使用闭环完成时延或梯度。Packet window 按 credit
payload 大小换算为在途 credit 数。\Cref{ssec:eval-cc} 在相同 credit gate 上
比较四个实例的利用率、队列和公平性。

\parab{DCQCN-style group ECN~\cite{dcqcn}.}
普通交换机按 egress queue 标记 CE。requester 把下行 credit/response 的 CE
回显到对应 normal request；上行路径保留沿途 CE，逐级聚合对所有输入取 OR。
aggregated packet 因而表示本轮任一分支是否被标记。responder 将标记样本比例
映射为 DCQCN 的拥塞程度，并按其降速与恢复状态机更新 $r_s$；$W_s$ 使用配置
的安全上限。CE 汇合复用 IP ECN 位和已有聚合头部状态，不引入 per-flow
交换机状态。

\parab{TIMELY~\cite{timely}.}
TIMELY 从连续闭环完成时延估计梯度，经平滑后更新 $r_s$；$W_s$ 使用配置的
安全上限。原算法中的 packet RTT 在此对应 \cref{eq:completion-delay} 的 $T_s$。
共享链路位于非关键分支时，梯度要到分支余量耗尽后才反映该链路的排队。

\parab{Swift~\cite{swift}.}
Swift 比较闭环完成时延与目标值，并按时延误差更新 $W_s$。窗口低于一份 credit
时，responder 设 $W_s=1$ 和 $r_s=P\cdot\mathrm{cwnd}/\mathrm{RTT}$，对应
$\mathrm{RTT}/\mathrm{cwnd}$ 的发放间隔。目标值由无队列完成时延加排队预算
得到。Swift 原有的 fabric/host 时延拆分在此合并为整条完成时延；单个样本无法
区分 fan-in 等待、requester 停顿和网络排队。

\parab{Oscar~\cite{oscar}.}
Oscar 同时使用时延、时延梯度、参考在途量和实际完成速率。Credit entry 在发出
时记录在途量，完成时把参考状态与 $T_s$ 交给批式估计器。控制律协调 window
loop 与 rate loop 后更新 $r_s$ 和 $W_s$；实际发放仍由
\cref{eq:cc-two-gates} 中更紧的门限制。该映射保留 Oscar 的参考状态同步，但
输入从单条点对点 RTT 变为 requester group 的最大完成时延，不能直接沿用其
单路径公平性结论。

具体控制律不属于 \sysname 的贡献。四个实例说明 rate/window gate 可以承载
显式标记和 completion delay 两类反馈。\Cref{ssec:eval-cc} 进一步量化
\cref{eq:branch-slack} 的分支过滤对利用率、公平性和队列的影响。

\section{真实队列上的分支过滤核对}
\label{app:cc-slack}

我们在 \texttt{asym-2348} 上保持四个成员、流量和有向 egress
$E=C3\rightarrow B2$ 不变，只改变 responder 放置。两个放置分别让经过
$E$ 的分支拥有余量，或使其成为当前最慢分支。独立无排队运行测得
$\Delta=12.664~\mu\mathrm{s}$；正式运行短时降低 $E$ 的服务率，并以 payload
packet 在 $E$ 上的 enqueue-to-service 中位数作为 $q_E$。每个放置包含一个
$q_E=0$ 对照和六个正 queue pulse；六个 pulse 在 $\Delta$ 两侧各取三个点。

\begin{figure*}[t]
\centering
\begin{tikzpicture}
\begin{axis}[
  arbor axis,arbor panel,axis on top,
  width=0.76\textwidth,height=0.34\textwidth,
  xmin=-0.6,xmax=24,ymin=-0.8,ymax=24,
  xtick={0,5,10,15,20},ytick={0,5,10,15,20},
  xlabel={Queueing delay added on $E$, $q_E$ ($\mu\mathrm{s}$)},
  ylabel={Delay visible in group completion ($\mu\mathrm{s}$)},
]
\path[fill=ArborSky,fill opacity=0.10,draw=none]
  (axis cs:0,-0.8) rectangle (axis cs:12.664,24);

% Parameter-free predictions frozen before the formal pulse runs.
\addplot[ArborBlue,thick,dashed,mark=none]
  coordinates {(0,0) (12.664,0) (24,11.336)};
\addplot[ArborOrange,thick,dash dot,mark=none]
  coordinates {(0,0) (24,24)};

% Measured medians from the two paired responder placements.
\addplot[ArborBlue,thick,only marks,mark=*,mark size=2.4pt,
  restrict expr to domain={\thisrow{critical}}{0:0}]
  table[col sep=comma,x=qUs,y=observedUs]
  {plots/data/pgf/cc_slack.csv};
\addplot[ArborOrange,thick,only marks,mark=square*,mark size=2.4pt,
  restrict expr to domain={\thisrow{critical}}{1:1}]
  table[col sep=comma,x=qUs,y=observedUs]
  {plots/data/pgf/cc_slack.csv};

\draw[ArborGray,densely dotted,line width=0.8pt]
  (axis cs:12.664,-0.8)--(axis cs:12.664,24);
\node[font=\scriptsize,text=ArborGray,anchor=north east]
  at (axis cs:12.35,23.4) {$\Delta=12.664~\mu\mathrm{s}$};

\node[font=\scriptsize,text=ArborBlue,align=left,anchor=west,
  fill=white,fill opacity=0.88,text opacity=1,inner sep=1.5pt]
  at (axis cs:14.2,5.0)
  {$E$ starts on a faster branch:\\only $\max(0,q_E-\Delta)$ is visible};
\node[font=\scriptsize,text=ArborOrange,align=left,anchor=west,
  fill=white,fill opacity=0.88,text opacity=1,inner sep=1.5pt]
  at (axis cs:10.4,17.6)
  {$E$ is on the slowest branch:\\all $q_E$ is visible};
\node[font=\scriptsize,text=ArborBlue,anchor=south]
  at (axis cs:6.3,0.7) {blue case, $q_E\le\Delta$: no completion change};

\node[font=\scriptsize,anchor=north west] at (axis cs:0.3,23.4)
  {markers: measured medians; lines: pre-declared predictions};
\end{axis}
\end{tikzpicture}

\caption{同一条 $E$ 队列是否进入整组完成反馈。$E$ 位于最慢分支时，每
$1~\mu\mathrm{s}$ 排队都增加约 $1~\mu\mathrm{s}$ 完成时延；$E$ 所在分支拥有
$\Delta=12.664~\mu\mathrm{s}$ 余量时，排队先消耗该余量，只有超出部分进入反馈。
标记为受 pulse 影响 offsets 的中位数；线为使用独立无排队运行所得
$\Delta$ 的无拟合参数预测。}
\label{fig:cc-slack}
\end{figure*}

\chapter{共享 allocator 的 slot 需求实验细节}
\label{app:state-depth}

\parab{拓扑与 workload.}
\Cref{ssec:eval-state} 使用 ls128 leaf-spine，并只报告同一台 focus spine
的状态。Host link 与 leaf--spine link 的单向速率均为 100~Gbps，传播时延为
3~\textmu s。每个 job 持续发送 8-GiB 的 3-rank 数据传输 workload，payload
为 8~KiB，使用一个 subchannel、4096 个共享 slots 和 32~MiB switch buffer。
job 数取 1、2、4、8、16、32、48 和 64；各 job 的起始时刻加入
$U(0,1~\mu\mathrm{s})$ 扰动。

共用端口配置运行 single-channel Broadcast，固定 sender host，并随 job
编号轮转其余两个成员。该放置保持一条共同的 sender-facing cut，但不要求
job 的完整 rank 集合相同。端口分散配置运行 per-rank AllGather；每个
AllGather 在状态核算中对应三个 3-rank Broadcast channels，spread placement
把不同 job 的 sender 分布到多个 leaf-facing 端口。这些端口共同推进同一
spine allocator。两类配置的 per-channel cold-start rate 分别为
1.5625 和 0.1953125~Gbps，并使用同一组冻结 Swift 参数：
startup window 为 4，target queue 为 5~\textmu s，flow-scaling range 为
10~\textmu s，cwnd 范围为 0.1--100，AI 为 0.5，$\beta$ 为 0.8，maximum
MDF 为 0.5，maximum pacing burst 为 1。它们用于实例化两种端口放置的容量
条件，不构成吞吐量基线比较。

\parab{所需 slot 数与理论界.}
主测量窗口为 $[150,190)$~ms，仿真运行到 200~ms 以排空 cohort。对该窗口中
分配的每个 normal context，仿真记录从本级 slot 绑定到本级收齐全部上行贡献
并完成聚合之间的 inclusive allocator sequence span。逐运行最大 span 即
exact cyclic reuse depth，也是该运行的实测所需 slot 数。聚合完成时结果已
形成，slot 进入可安全覆盖状态；旧内容可保留到 allocator 再次访问时。

原始记录还区分推进 allocation 的 packet payload 长度。分析在聚合前要求
focus spine 的 \texttt{payloadSizedAllocs} 等于全部 \texttt{allocs}，且
\texttt{headerOnlyAllocs}=0；任一 header-only credit 推进 allocator 都使
该运行失败。该检查核对
\cref{eq:allocator-payload-carrier-bound} 的 carrier 前提。分析再对
\texttt{run-switch\_perchan.csv} 中测量窗口边界的累计 allocation 计数作差，
并与 \texttt{run-messages.csv} 的 responder 放置连接，得到该次 cohort
实际参与推进的有向端口集合 $\mathcal P_e^{\rm obs}$。逐 channel 增量之和
必须等于 focus spine 的总 allocation 数。

令 $R_{\max}^{\rm obs}$ 为同一运行中 normal credit 的最大端到端 completion
delay。由于本级 slot 区间包含在该闭环内，用
$R_{\max}^{\rm obs}$ 代替本级生命周期得到保守理论界。令 $P=8192$ bytes，
分析逐运行计算
\begin{align}
A_{\rm common}^{\rm obs}
  &=\left\lceil\frac{C_{p^\star}R_{\max}^{\rm obs}}{8P}\right\rceil,\\
A_{\mathcal P}^{\rm obs}
  &=\sum_{p\in\mathcal P_e^{\rm obs}}
    \left\lceil\frac{C_pR_{\max}^{\rm obs}}{8P}\right\rceil .
\end{align}
共用端口配置要求实测所需 slot 数不超过
$A_{\rm common}^{\rm obs}$；端口分散配置要求其不超过
$A_{\mathcal P}^{\rm obs}$。逐端口取 ceiling 保留并行 packet 的边界效应。
实验代入 $R_{\max}^{\rm obs}$ 核对两条关系；部署时应代入
部署 contract 保证的 $R_{\max}$。

\parab{重复、无漂移检查与完整性.}
每个点运行五个 ns-3 RNG run，并报告双侧 Student-t 95\% CI（$df=4$）。
每次运行至少要求 10,000 个 normal-credit 样本。16、32 和 64 jobs 另在
$[350,390)$~ms 复测，并运行到 400~ms。
预先规定最大 normal-credit completion delay、实测所需 slot 数和 allocator
rate 的早晚窗口五次运行均值变化均不超过 10\%。

完整矩阵包含 110 次运行；最小 normal-credit cohort 有 58,261 个样本。
共用端口放置的 55 次运行全部满足单端口理论界，最小逐运行余量为 25 slots；
端口分散放置的 55 次运行全部满足实际参与端口集合的求和界，最小余量为
107 slots。最大 completion delay、实测所需 slot 数和 allocator rate 的
早晚窗口均值变化最大为 8.89\%。

分析没有纳入失败、删失或选择性排除的运行。每次运行的 16 个 switch cohort
均在停止前排空，focus spine 的 allocation 与本级完成计数一致。全部
focus-spine allocation 均由至少携带 8192-byte payload 的 response 推进，
header-only allocation 为零；所有运行还同时满足无 PFC pause、
\texttt{AGG\_MISS}、repair、提前覆盖和协议违例的 clean-run 条件。

\chapter{固定基础设施多租户实验细节}
\label{app:state-capacity}

\parab{固定拓扑与放置.}
\Cref{fig:multitenant-memory} 固定一台 32$\times$200-Gbps 被测 spine，
单链路传播时延为 3~\textmu s。内存扫描的 Arbor/SwitchML-like 拓扑装配
使用四个 access domain，每个 domain
以一条 1.6-Tbps bundle 表示八个物理端口；每个八-rank job 以 $4+4$ 放在
一对 domain 上，job 交替使用两对 domain，$J=12$ 时每台 leaf 使用 24 张
互不复用的 NIC。ATP-like 使用其贯穿全文的 worker-ToR/PS-ToR U200
装配；其 worker access 与 fabric link 均为 200~Gbps，但每 job
logical offer 固定为 100~Gbps，该 access link 不比 bundle 装配更紧。

\Cref{fig:multitenant-population} 使用总端口容量相同的 64$\times$100-Gbps
spine；仿真将四条物理 lane 合并为一条 400-Gbps access-domain bundle。每个
八-rank job 放在一对 domain 上，两侧各有四个独占 NIC；job 增加后只竞争
domain--spine carrier。两组实验的 leaf 聚合状态均充足，被测状态池只位于
root。人口扫描按各自的拓扑预算归一化，不跨拓扑比较绝对 Gbps。

\parab{p99 operating point 的取得.}
我们先在人口扫描的固定 6.4-Tbps spine、256-MiB AllReduce 和充足状态下运行
三次，
在 $[40,100)$~ms 稳定窗口记录 normal completion delay。三次精确样本数为
1,041,345、1,046,260 和 1,042,882，p99 分别为 51,906、49,692 和
48,421~ns，最大值分别为 101,103、101,529 和 100,348~ns；所有 cohort 均在
停止前排空，且无 \texttt{AGG\_MISS}、repair、PFC 或丢包。
我们将最坏 p99 向上取整为 $R_{99}=55$~\textmu s，并按
\cref{eq:allocator-payload-carrier-bound} 配置
\begin{equation}
  64\left\lceil\frac{100\,\mathrm{Gbps}\times55\,\mu\mathrm{s}}
  {8\times8192\,\mathrm{bytes}}\right\rceil
  =64\times84=5376\ \text{slots}=42\ \mathrm{MiB}.
\end{equation}
该数值是实测 p99 operating point。因为 pilot 仍有超过 55~\textmu s 的 tail，
它不构成部署 hard guarantee；需要 hard guarantee 时仍须使用
\cref{ssec:bounded-state} 所要求的契约 $R_{\max}$。

\parab{内存扫描的 workload 与状态变量.}
\Cref{fig:multitenant-memory} 中，每个 job 连续执行 16 轮 256-MiB
AllReduce，不插入 compute gap；iteration~0 用于 warm-up，rank submit
skew 为 8~\textmu s，
compute 与 start jitter 均为 0。这样 $J$ 就是实测窗口内同时通信的 job 数，
而不是由 communication duty 推断出的名义并发。每个配置运行三次。

spine 总内存逐 MiB 取 1--8~MiB，即 128--1024 个 8-KiB slots。
SwitchML-like 始终按配置人口 $K=16$ 静态等分总 slots；每个 logical index
由两个交替使用的 physical version sets 支撑，因此 logical window 是分区
深度的一半。每 MiB 总内存为每个配置 job 提供八个 physical slots 和四个
logical indices。正文图统一显示 2--8~MiB；未显示的 1-MiB 原始点仍保留在
正式 CSV 中。该点为 SwitchML-like 提供 4 个 active slots/job，goodput 为
19.55~Gbps，且 hash collision 为零；host window 在 root 插入前阻止越界
context，因此静态分配把状态不足转化为端侧 window stall，而不是 root
聚合冲突。\sysname 使用 switch-wide cyclic
allocator。ATP-like 使用
cumulative expected sequence，开启 collision-driven remap、PS fallback
及其 ECN/AIMD 路径；初始与最大窗口为 128/512 个 8-KiB packet，每个 job
配置 512 个逻辑 hash entries。每个 job 使用独立 100-Gbps PS access link 和
100-Gbps fallback reducer；normal aggregate 与 collision fallback 共用该链路。
两个 baseline 均采用各自的 idealized 8-KiB payload 语义，用于隔离状态管理
而非复刻生产实现的全部开销。
本实验中的 AllReduce 两个方向均携带 payload；spine allocator 的 normal
advance 均由携带 payload 的 response 推进。

\parab{内存扫描的测量与完整性.}
对每次运行，分析器以所有 job 的 iteration~0 完成时间最大值作为公共窗口起点，
以所有 job 的 iteration~15 完成时间最小值作为终点。仅当 collective 中点落在
该窗口内时才计入，并逐点验证该中点上恰有全部 $J$ 个 job 在通信，且每个 job
至少有四个样本。每个 collective 的 logical goodput 定义为
$256\,\mathrm{MiB}\times8/\mathrm{CCT}$；分析先在每个 job 内取均值，再对
job 等权平均，最后跨三次运行报告均值与双侧 Student-t 95\% CI（$df=2$）。
图中的百分比仅将 Gbps 值除以固定 100~Gbps，不随观测并发改变分母。

完整矩阵包含 $3\times8\times3\times3=216$ 个 task；全部 task 完成，
公共饱和窗口共保留 25,557 个 collectives。没有失败 task 或选择性剔除的
重复运行。全部运行在停止前排空 leaf 和 spine 状态，无协议违例或 PFC
pause，充足 leaf 层始终为零 \texttt{AGG\_MISS}。稀缺 spine 点中的
\texttt{AGG\_MISS}、repair 和覆盖均保留在结果中，而非作为异常删去。
72 个 SwitchML-like task 均实际使用两个 physical version sets，并在停止时
清空聚合状态。72 个 ATP-like task 均执行 cumulative expected-sequence transport；
高序 parameter packet 不释放发送窗口，缺口由快速重传或 timeout 修复。ATP-like
总计产生 49,152,582 次 root collision bypass 和 98,324,256 个 PS fallback
packets；这些 packet 全部经过每 job 独立的 100-Gbps PS 链路与 reducer。

\begin{figure*}[t]
\centering
\begin{tikzpicture}
\begin{groupplot}[
  group style={group size=3 by 1,horizontal sep=0.025\textwidth},
  arbor axis,
  arbor panel,
  width=0.315\textwidth,height=0.30\textwidth,
  xmin=1.7,xmax=8.3,ymin=0,ymax=30,
  restrict x to domain=2:8,
  xtick={2,3,4,5,6,7,8},
  ytick={0,10,20,30},
  xticklabel style={font=\scriptsize},
  xlabel={Spine memory (MiB)},
]
\nextgroupplot[
  title={(a) $J=4$},
  ylabel={Shared-map conflicts / attempts (\%)},
  legend to name=multitenantconflictlegend,
  legend columns=2,
  legend style={/tikz/every even column/.append style={column sep=5pt}},
]
\addplot[ArborBlue,thick,mark=*,error bars/.cd,y dir=both,y explicit]
  table[col sep=comma,x=memoryMiB,y=arbor,y error=arborCi,
    restrict expr to domain={\thisrow{jobs}}{4:4}]
    {plots/data/pgf/multitenant-capacity-conflicts.csv};
\addlegendentry{Arbor}
\addplot[ArborOrange,thick,dashed,mark=square*,
  error bars/.cd,y dir=both,y explicit]
  table[col sep=comma,x=memoryMiB,y=atp,y error=atpCi,
    restrict expr to domain={\thisrow{jobs}}{4:4}]
    {plots/data/pgf/multitenant-capacity-conflicts.csv};
\addlegendentry{ATP-like}

\nextgroupplot[
  title={(b) $J=8$},
  yticklabels={},
]
\addplot[ArborBlue,thick,mark=*,error bars/.cd,y dir=both,y explicit]
  table[col sep=comma,x=memoryMiB,y=arbor,y error=arborCi,
    restrict expr to domain={\thisrow{jobs}}{8:8}]
    {plots/data/pgf/multitenant-capacity-conflicts.csv};
\addplot[ArborOrange,thick,dashed,mark=square*,
  error bars/.cd,y dir=both,y explicit]
  table[col sep=comma,x=memoryMiB,y=atp,y error=atpCi,
    restrict expr to domain={\thisrow{jobs}}{8:8}]
    {plots/data/pgf/multitenant-capacity-conflicts.csv};

\nextgroupplot[
  title={(c) $J=12$},
  yticklabels={},
]
\addplot[ArborBlue,thick,mark=*,error bars/.cd,y dir=both,y explicit]
  table[col sep=comma,x=memoryMiB,y=arbor,y error=arborCi,
    restrict expr to domain={\thisrow{jobs}}{12:12}]
    {plots/data/pgf/multitenant-capacity-conflicts.csv};
\addplot[ArborOrange,thick,dashed,mark=square*,
  error bars/.cd,y dir=both,y explicit]
  table[col sep=comma,x=memoryMiB,y=atp,y error=atpCi,
    restrict expr to domain={\thisrow{jobs}}{12:12}]
    {plots/data/pgf/multitenant-capacity-conflicts.csv};
\end{groupplot}
\node[anchor=south,yshift=26pt] at (group c2r1.north)
  {\ref{multitenantconflictlegend}};
\end{tikzpicture}

\caption{内存扫描中的共享映射冲突诊断。Arbor 的事件为 root
\texttt{AGG\_MISS}，ATP-like 的事件为 root collision bypass；分母均为
root aggregation attempts。三个面板使用与\cref{fig:multitenant-memory}
相同的 $J$、2--8~MiB 内存横轴和三次运行。SwitchML-like 不执行共享地址
映射：静态 active bank 耗尽时由 host window 在 root 插入前停发，因此不画
一条无信息量的零冲突曲线。标记为三次运行的均值，误差线为双侧 Student-t
95\% CI。}
\label{fig:multitenant-pressure}
\end{figure*}

该诊断不能被解释为三种方案的统一性能指标。以 $J=12$ 为例，Arbor 在
2、4 和 8~MiB 的 \texttt{AGG\_MISS}/attempt 分别为 3.92\%、0.26\% 和零；
ATP-like 在相同三点的 collision/attempt 为 16.04\%、8.97\% 和 5.73\%，
对应 fallback context rate 为 27.64\%、16.46\% 和 10.84\%。SwitchML-like
的状态压力则表现为正文所述的 active-bank window stall；它没有 shared-map
collision，不表示其静态分区仍有可供活跃 job 使用的状态。

\parab{约 75\% 活跃人口的 configured-population 扫描.}
\Cref{fig:multitenant-population} 固定 42~MiB spine 状态并扫描
\[
K\in\{16,32,64,128,192,256,336,512\}.
\]
每个点约 75\% 的 job 持续通信，其余 job 在该快照中保持静默。实际活跃人口和
启动偏移记录在 paper-data CSV 与 immutable manifest 中。每个活跃 job 连续
执行四轮 256-MiB AllReduce，iteration~0 为 warm-up；三个确定性运行将相同的
启动偏移逐项复用于三种方案。

令 $n$ 为最繁忙 domain pair 上的活跃 job 数。ATP-like 与 SwitchML-like 的
per-job network-only service budget 为
$\min(100,400/n)$~Gbps；Arbor 的 $4+4$ AllReduce 在该 carrier 上具有实测的
1.5 倍 wire factor，其预算为 $\min(100,400/(1.5n))$~Gbps。发送端加载冻结
预算的 98\%，图中仍除以完整预算。该分母只由拓扑、放置和活跃人口计算，不由
观测吞吐反推。

SwitchML-like 为全部 $K$ 个 job 静态预留状态；每-job 分区依次为
336、168、84、42、28、21、16 和 10 个 physical slots。每份分区再均分为两个
交替的 physical version banks，因此 active bank 依次含
168、84、42、21、14、10、8 和 5 slots。Arbor 和 ATP-like 不为未活跃 job
分配运行时状态。三种方案均对每个 $K$ 运行三次，共 72 个 task。

分析只保留该点全部活跃 job 均在通信的公共饱和窗口，先在 run 内对 job
等权，再对三次运行报告双侧 Student-t 95\% CI（$df=2$）。72 个 task
全部通过 completion、placement、root-memory、lifecycle 和 PFC 审计。SwitchML-like
的 alternating banks 在停止时均清空；Arbor 的 spine \texttt{AGG\_MISS}、repair
和 PFC 均为零。ATP-like 的 collision/preaggregation-output 比例随 $K$ 从
0.007\% 增至 16.87\%；该模型的冲突 offset 转入每-job PS fallback 路径。

\parab{暂时离开干预.}
\Cref{fig:multitenant-state} 的
动态离开实验使用 U200 topology、608-KiB root 状态、3 个 endpoint-disjoint
job 和 8 ranks/job。每个 job 不插入 compute gap，连续执行 48 轮
256-MiB AllReduce；前 23 轮把干预移出启动阶段。第三个 job
完成 iteration~23 后额外等待 250~ms，再提交 iteration~24；另外两个 job
的调度不变。三次冻结运行使用配对的初始启动偏移；ATP-like 的 hash seed
随运行编号独立变化。分析器从原始完成时间定位各方案的暂停边界，并将其对齐为
$t=0$；job 在 $t=250$~ms 返回。collective 按完成时刻进入 50-ms bin，iteration~0
不计。每个 bin 先在 job 内求均值，再对两个未暂停 job 等权；正文曲线为三次
运行的均值，阴影为最小值至最大值。

正文对称报告 $[-100,350]$~ms，即离开前和返回后各 100~ms；该区间内三个
方案、三次运行和两个未暂停 job 在每个 bin 均有样本。350--400~ms bin 中至少
一个未暂停 job 没有 collective 完成，因此分析器不延续最后一次观测，也不据此
判断返回后的长期收敛。响应时间单独由事件后提交的首个
完整 collective 计算。9 个 task 均通过 completion、placement、单次提交守恒和
PFC 审计。Arbor 的 allocator-pressure profile 允许 \texttt{AGG\_MISS} 和 repair；
停止时 root state 为空，且没有协议违例。

\parab{初始启动安排的敏感性.}
我们另在冻结的 $J=8,M=4$~MiB workload 上比较 job 同时启动、
均匀错开和五组预先生成的启动安排。不同安排改变 ATP-like 的数值，
但没有改变三种方案的排序；因此正文各图不把它作为扫描变量。

\begin{table}[t]
\centering
\small
\begin{tabular}{lrrrr}
\toprule
Scheme & Aligned & Staggered & Random range & Max dev. \\
\midrule
Arbor & 99.455 & 99.455 & 99.455--99.455 & 0.00\% \\
ATP-like & 80.059 & 95.195 & 88.691--94.006 & 11.37\% \\
SwitchML-like & 78.532 & 78.532 & 78.532--78.532 & 0.00\% \\
\bottomrule
\end{tabular}
\caption{不同初始启动安排下的 per-collective logical goodput (Gbps)。
Random range 覆盖五组预先生成的启动安排。}
\label{tab:g5-phase}
\end{table}
